% 【本文件写什么】impl 实现层：fs-bridge
%   - 定位：VFS 与 wateros-fs 桥接后端。
%   - crate 路径：os/components/wateros-vfs/vfs-impl/impl-fs-bridge/src/
%   - 如何实现对应 api-v0 trait：关键类型、状态机、数据结构
%   - 算法或硬件交互流程（可用伪代码/流程图）
%   - 本 impl 启用的 Cargo [features] 与 arch feature（impl-riscv64 等）
%   - 与其它 impl 变体的差异、切换 feature 时的影响
%   - 性能/内存假设与已知限制
% 【事实来源】
%   - 上述 crate 源码与 Cargo.toml
%   - os/feature-tree.txt 中指向本 impl 的 feature 链

\subsubsection{impl-fs-bridge}

\paragraph{总体结构。}
零大小类型\code{FsBridge}实现\code{VfsBackend}，将路径级操作委托给\code{wateros-fs}的\code{rootfs}、\code{devfs}、\code{procfs}与ext4卷trait，不re-export \code{FsError}等FS类型。子模块：\code{mount\_table}（路由解析）、\code{mount\_ns}（辅助挂载）、\code{paged\_handle}（页缓存句柄）、\code{file\_handle}/\code{dir\_handle}/\code{proc\_handle}、\code{tmpfs}（内存伪卷）。

\paragraph{路径路由。}
\code{resolve\_route(path)}返回\code{FsRoute}枚举：根卷\code{Root}、辅助RO/RW卷\code{AuxRo}/\code{AuxRw}、\code{PseudoProc}、\code{PseudoSecurity}等。\code{metadata}/\code{read}/\code{open}按路由选择\code{SharedRwFs}、\code{proc\_view()}或tmpfs状态机。普通文件\code{metadata}对根卷文件叠加页缓存\code{logical\_size}。

\paragraph{打开与I/O。}
\code{open}根据节点类型返回\code{PagedFileHandle}（\code{FILE\_IO\_MODE=Direct}时走页缓存）、\code{DirectoryHandle}、\code{ProcHandle}或字符设备句柄（转交fd-session）。\code{PagedFileHandle}经\code{FsPageIo}实现\code{PageCacheIo}，miss/flush时短持\code{SharedRwFs}锁调用\code{write\_range}/\code{read\_range}。

\paragraph{挂载命名空间。}
对外导出（经聚合层转发）：\code{mount\_ext4\_block\_at}、\code{mount\_tmpfs\_at}、\code{mount\_cgroup\_at}、\code{mount\_procfs\_at}、\code{mount\_securityfs\_at}、\code{mount\_bind\_at}、\code{move\_mount\_at}、\code{set\_mount\_propagation}、\code{remount\_readonly\_at}、\code{unmount\_at}（支持lazy detach）、\code{assert\_path\_writable}、\code{mount\_statfs\_magic}。\code{mount\_procfs\_at}前注册procfs回调并确保\code{/proc}挂载点存在。

\paragraph{路径级元数据写操作。}
\code{mkdir\_path}、\code{unlink\_path}、\code{rename\_path}、\code{chmod\_path}、\code{chown\_path}、xattr族、\code{truncate\_path}、\code{symlink\_path}、\code{mknod\_socket\_at}、\code{overwrite\_file\_at}（unlink+写+页缓存驱逐）均先\code{resolve\_route}再下探对应卷或伪FS。

\paragraph{错误与代次。}
\code{map\_fs\_err}集中\code{FsError}$\rightarrow$\code{VfsError}。辅助卷挂载/卸载与\code{fs::rootfs::mount\_generation}联动；\code{reset\_file\_page\_cache}批量刷回并回收全局页缓存。

\paragraph{限制。}
securityfs/cgroup/tmpfs为满足LTP子集的伪实现；bind/move传播类型覆盖有限。页缓存与ext4写路径锁顺序须严格遵守（见\code{page-cache}与\code{paged\_handle}文档），否则可能死锁。
